Autonomous tool-using Large Language Model (LLM) agents are increasingly deployed in complex domain-specific workflows. However, deploying agents in healthcare-commerce applications presents critical security and policy challenges, including unauthorized order modifications, prescription bypasses, cross-user data exposure, and indirect prompt injections. In this paper, we present \textbf{MediRush-SafeAgent}, a formal multi-stage runtime framework designed to guarantee policy-constrained, authorized, and reliable tool execution in healthcare-commerce workflows. MediRush-SafeAgent integrates a pre-execution domain policy engine, a strict scope authorization verifier, a post-execution environment state verifier, and dynamic risk-adaptive human escalation. We construct \textbf{MediRushBench}, a standardized evaluation benchmark comprising 24 scenarios across 12 operational and security categories. Our empirical evaluation demonstrates that MediRush-SafeAgent achieves a 100.0\% Safe Task Completion Rate ($STCR$) and 0.0\% Policy Violation Rate ($PVR$), outperforming unconstrained LLM baselines ($STCR = 83.33\%$, $UTIR = 16.67\%$) and prompt-guardrail baselines ($STCR = 95.83\%$). These results establish the efficacy of joint multi-stage runtime interception for domain-constrained agent safety.
